\documentclass[11pt]{article}

\usepackage[margin=1in]{geometry}
\setlength{\headheight}{14pt}
\usepackage[T1]{fontenc}
\usepackage[utf8]{inputenc}
\usepackage{xcolor}
\usepackage{hyperref}
\usepackage{booktabs}
\usepackage{longtable}
\usepackage{enumitem}
\usepackage{tikz}
\usetikzlibrary{arrows.meta, positioning, shapes.multipart}
\usepackage{titling}
\usepackage{fancyhdr}

\definecolor{linkblue}{HTML}{1155CC}
\hypersetup{
    colorlinks=true,
    linkcolor=linkblue,
    urlcolor=linkblue,
    citecolor=linkblue,
    pdftitle={Offline Real-Time Voice Assistant -- Submission},
    pdfauthor={Vichruth M}
}

\pagestyle{fancy}
\fancyhf{}
\rhead{Offline Real-Time Voice Assistant}
\lhead{AI Engineering Intern Assignment}
\cfoot{\thepage}
\renewcommand{\headrulewidth}{0.4pt}

\setlength{\droptitle}{-2cm}
\emergencystretch=1.5em

\title{\textbf{Offline Real-Time Voice Assistant}\\[4pt]
\large Design Document, Measured Latency, AI-Usage Disclosure \& Deliverable Links}
\author{Vichruth M\\ \href{mailto:vichruth.m2024@vitstudent.ac.in}{vichruth.m2024@vitstudent.ac.in}}
\date{15 July 2026}

\begin{document}
\maketitle
\thispagestyle{fancy}

\section*{Objective}
Build a real-time, audio-in / audio-out conversational AI assistant that accepts voice queries
and responds with voice, targeting an end-to-end response time within 2 seconds wherever
possible, with a fallback flow that keeps the user engaged through delays or interruptions
instead of failing silently or with a generic error. This submission implements the assistant
fully offline: speech recognition, language generation, and speech synthesis all run locally,
with no network calls at any stage.

\section{Architecture}

\begin{center}
\resizebox{\textwidth}{!}{%
\begin{tikzpicture}[
    node distance=6mm and 8mm,
    box/.style={draw, rounded corners, minimum height=9mm, minimum width=20mm, align=center, font=\small},
    lbl/.style={font=\tiny, align=center, text width=22mm}
]
\node[box] (mic) {Mic};
\node[box, right=of mic] (vad) {Silero\\VAD};
\node[box, right=of vad] (ack) {Filler\\ack};
\node[box, right=of ack] (stt) {faster-\\whisper};
\node[box, right=of stt] (llm) {qwen2.5:3b\\(Ollama)};
\node[box, right=of llm] (tts) {Piper\\TTS};
\node[box, right=of tts] (spk) {Speaker};

\draw[-Stealth] (mic) -- (vad);
\draw[-Stealth] (vad) -- (ack) node[midway, above, font=\tiny] {endpoint};
\draw[-Stealth] (ack) -- (stt);
\draw[-Stealth] (stt) -- (llm) node[midway, above, font=\tiny] {transcript};
\draw[-Stealth] (llm) -- (tts) node[midway, above, font=\tiny] {sentence};
\draw[-Stealth] (tts) -- (spk);

\node[lbl, below=2mm of vad] {endpointing\\$<$1 ms/frame, CPU};
\node[lbl, below=2mm of ack] {$<$300 ms\\pre-rendered};
\node[lbl, below=2mm of stt] {STT, GPU\\small.en, int8};
\node[lbl, below=3mm of llm] {LLM, streamed\\3B params};
\node[lbl, below=2mm of tts] {TTS, CPU\\streamed};
\end{tikzpicture}%
}
\end{center}

\noindent\textbf{The core latency trick.} A naive pipeline waits for the complete LLM response
before synthesizing and playing audio -- typically 4--6 seconds for a short reply, which fails
the 2-second brief. Instead, tokens are streamed from the LLM and split on sentence boundaries
(\texttt{.\ ! ? \textbackslash n}, or a comma once a clause exceeds \textasciitilde12 words).
The instant the first complete sentence is available it is sent to Piper and played, while the
LLM continues generating the remaining sentences in the background. The metric that matters is
\textbf{time-to-first-audio (TTFA)}, not total generation time, since that is what the user
actually perceives as latency.

\subsection*{Design decisions}
\begin{center}
\begin{longtable}{@{}p{2.1cm}p{3.1cm}p{8.3cm}@{}}
\toprule
\textbf{Component} & \textbf{Choice} & \textbf{Rationale} \\
\midrule
VAD & Silero VAD & Neural VAD, far less prone to false triggers on fan/keyboard noise than
WebRTC or energy-based detectors; cheap enough to run per 32\,ms frame on CPU, leaving the
GPU free. \\
STT & faster-whisper, \texttt{small.en}, int8\_float16, CUDA & CTranslate2 backend runs
\textasciitilde4$\times$ faster than reference Whisper. \texttt{base.en} loses technical
vocabulary; \texttt{medium} is too slow for the latency budget on a 6\,GB GPU.
\texttt{small.en} is the accuracy/latency knee for this hardware. \\
LLM & \texttt{qwen2.5:3b} via Ollama, streaming & A 7B model does not fit the 2\,s budget on
a 6\,GB laptop GPU; a 3B model gives a first token in \textasciitilde180--230\,ms. Deliberate
latency-vs-quality tradeoff, mitigated by sentence-level streaming so the user hears audio
long before generation finishes. \\
TTS & Piper TTS, CPU (voice \texttt{lessac}\--\texttt{medium}) & Lightweight VITS model, \textasciitilde50--100\,ms to synthesize a sentence. Deliberately kept on CPU so the 6\,GB GPU is shared only
between STT and the LLM, avoiding VRAM contention. Coqui/XTTS sound better but take seconds
per utterance and are not viable inside the budget. \\
Audio I/O & \texttt{sounddevice} (PortAudio) & Callback-based streaming I/O, which is what
makes low-latency playback and immediate barge-in cancellation possible. \\
\bottomrule
\end{longtable}
\end{center}

GPU memory with STT, LLM, and TTS models all resident: \textasciitilde3.6\,GB of 6\,GB
available (RTX 4050 laptop GPU).

\section{Fallback State Machine}

The assistant is explicitly designed to never stall silently, never show a generic error, and
never end a conversation abruptly. This is treated as a first-class state machine, not an
afterthought:

\begin{center}
\begin{tikzpicture}[
    node distance=14mm and 20mm,
    state/.style={draw, rounded corners, minimum width=26mm, minimum height=9mm, align=center, font=\small},
    >=Stealth
]
\node[state] (idle) {IDLE};
\node[state, right=of idle] (listen) {LISTENING};
\node[state, right=of listen] (think) {THINKING};
\node[state, right=of think] (speak) {SPEAKING};

\draw[->] (idle) -- (listen) node[midway, above, font=\tiny] {startup};
\draw[->] (listen) -- (think) node[midway, above, font=\tiny, align=center] {utterance\\endpointed};
\draw[->] (think) -- (speak) node[midway, above, font=\tiny, align=center] {first sentence\\synthesized};
\draw[->] (speak) to[bend left=40] node[midway, above, font=\tiny] {barge-in detected} (listen);
\draw[->] (speak) to[bend left=55, looseness=1.4] node[midway, below=1mm, font=\tiny] {reply complete} (listen);
\end{tikzpicture}
\end{center}

\begin{description}[leftmargin=4mm, itemsep=3pt]
\item[Layer 1 -- instant acknowledgement ($<$300\,ms).] Eight short filler lines (``Mhm.'',
``Let me think.'', ``Good question.'', ...) are synthesized once at startup and cached in
memory. The moment VAD detects end-of-speech, a random filler plays immediately, in parallel
with STT and LLM startup, so there is never a silent gap while the pipeline warms up.
\item[Layer 2 -- watchdog escalation.] If no LLM token has arrived after 1.5\,s, a contextual
stall line plays (``that's an interesting one, give me a second''). If nothing has arrived by
5\,s, the turn is abandoned: an honest holding reply plays and the microphone reopens. The
assistant never states that an error occurred and never terminates the session. If Ollama is
unreachable at all -- including if it is killed mid-conversation -- the same degrade path
fires: the assistant acknowledges, apologizes briefly, and invites the user to continue.
\item[Layer 3 -- barge-in.] If the user speaks while the assistant is talking, sustained
speech probability above threshold for \textasciitilde200\,ms flushes the playback queue,
stops the audio stream, and cancels the in-flight LLM generation via a cooperative cancel
event, returning immediately to LISTENING. A short grace window after playback starts and a
higher confirmation threshold than normal endpointing prevent the assistant's own voice from
re-triggering itself on speaker leakage.
\end{description}

\section{Measured Latency}

Latency is timestamped per stage relative to end-of-speech and logged to
\texttt{latency\_log.csv} on every turn (see \texttt{metrics.py}); \texttt{test\_turn.py} runs
the pipeline headlessly on fixed transcripts to reproduce these numbers without a microphone.

\begin{center}
\begin{tabular}{@{}lr@{}}
\toprule
\textbf{Stage} & \textbf{Median} \\
\midrule
Filler acknowledgement & \textasciitilde90\,ms \\
STT (faster-whisper small.en) & \textasciitilde200\,ms \\
LLM first token & \textasciitilde180--230\,ms \\
\textbf{Time to first audio (TTFA)} & \textbf{\textasciitilde520--600\,ms} \\
\bottomrule
\end{tabular}
\end{center}

\begin{quote}\small
\color{red!70!black}
\textbf{Before submitting:} regenerate this table from a fresh \texttt{latency\_log.csv} (run
\texttt{python test\_turn.py} or a live session) and paste the actual medians here. Do not
submit invented numbers -- reviewers may ask for the raw CSV.
\end{quote}

\section{AI-Usage Disclosure}

Anthropic's Claude (Claude Code) was used as a development assistant throughout this project.
Specifics:

\begin{itemize}[leftmargin=4mm]
\item \textbf{Scaffolding \& implementation.} Claude Code generated the initial implementation
of \texttt{main.py}, \texttt{vad.py}, \texttt{stt.py}, \texttt{llm.py}, \texttt{tts.py},
\texttt{state.py}, \texttt{audio\_io.py}, and \texttt{metrics.py} from a written specification
(\texttt{context.md}) describing the required pipeline, latency budget, and fallback behaviour.
\item \textbf{What it got wrong / had to be corrected.}
\begin{quote}\small\color{blue!40!black}
REPLACE WITH SPECIFICS -- e.g. the exact fix needed for the ALSA/PipeWire device conflict in
\texttt{run.sh} (the onboard codec's analog profile has to be released from PipeWire before
raw ALSA access works, which the first draft didn't handle), any tuning of
\texttt{BARGE\_THRESHOLD}/\texttt{BARGE\_CONFIRM} needed after testing on laptop speakers
(speaker-to-mic leakage triggering false barge-in), or watchdog timer values changed after
observing real generation latency. Name the file and the concrete before/after.
\end{quote}
\item \textbf{Design doc \& this submission document.} This PDF's structure and the README were
drafted with Claude Code assistance and reviewed/edited for accuracy against the actual
measured behaviour of the code before submission.
\item \textbf{What was not AI-generated.} REPLACE -- e.g. hardware-specific tuning done by
running on your own machine, any manual verification of the measured latency numbers, the demo
video, and final review of every file.
\end{itemize}

\section{Deliverables}

\begin{center}
\begin{tabular}{@{}p{4.2cm}p{9.3cm}@{}}
\toprule
\textbf{Deliverable} & \textbf{Link} \\
\midrule
GitHub repository & \url{REPLACE-WITH-GITHUB-URL} \\
Repository zip (Google Drive) & \url{REPLACE-WITH-DRIVE-LINK} \\
Demo video (Drive) -- shows on-screen TTFA, Ollama killed mid-conversation with graceful
degradation, and barge-in working & \url{REPLACE-WITH-DRIVE-LINK} \\
This design document (PDF, Drive) & \url{REPLACE-WITH-DRIVE-LINK} \\
\bottomrule
\end{tabular}
\end{center}

\begin{quote}\small\color{red!70!black}
Verify every link above in an incognito / logged-out browser window before submitting --
broken or permission-restricted links are treated as an automatic reject.
\end{quote}

\end{document}
